\frontmatter												%Seitennumerierung
\pagenumbering{Roman}										%Römische Zahlen
\addtocounter{page}{2}

\newcommand{\doublesignature}[2]{%
  \parbox{\textwidth}{
    \hfill
    \parbox{7cm}{
      \centering
      \rule{6cm}{1pt}\\
      #1
    }
    \parbox{7cm}{
      \centering
      \rule{6cm}{1pt}\\
      #2
    }
  }
  \mbox{}\\
  \mbox{}\\
  \mbox{}\\
  \mbox{}\\
}

\vspace*{20pt}

\section*{Eidestattliche Erklärung}
\label{sec:eidestattliche-erklaerung}
Ich erkläre an Eides statt, dass ich die vorliegende Arbeit selbstständig verfasst, andere als die angegebenen
Quellen/Hilfsmittel nicht benutzt und die den benutzten Quellen wörtlich und inhaltlich entnommenen
Stellen als solche kenntlich gemacht habe.\\
\\
Arnfels, am 05. April 2018\\

\vskip 2cm

\doublesignature{Lukas Freyler}{Stefan Ornik}
\doublesignature{Fabio Pölzl}{Dominik Riegelnegg}

\vskip 5cm

\clearpage

\newpage
\thispagestyle{empty}
\mbox{}

\clearpage

\section*{Danksagung}
\label{sec:danksagung}
An dieser Stelle möchten wir uns bei allen bedanken, die uns im Rahmen der Diplomarbeit unterstützt und betreut haben.\\
\\
Besonders bedanken möchten wir uns bei unseren Betreuern, Dipl.-Ing. Werner Harnisch, Dipl.-Ing. Wolfgang Mader, BEd. Otto Schuller
und Dipl.-Ing. Manfred Steiner für die fachliche Unterstützung während dieser Arbeit.

\clearpage

\newpage
\thispagestyle{empty}
\mbox{}

\clearpage

\section*{Abstract}
\label{sec:abstract}
TODO

\section*{Zusammenfassung}
TODO

\clearpage

\newpage
\thispagestyle{empty}
\mbox{}

\clearpage

\subsection*{Gender Erklärung}
\label{sec:gender-erklaerung}
Aus Gründen der besseren Lesbarkeit wird in dieser Arbeit die Sprachform des generischen Maskulinums angewendet. Es wird an dieser Stelle darauf hingewiesen, dass die ausschließliche Verwendung der männlichen Form geschlechtsunabhängig verstanden werden soll.

\subsection*{Über dieses Dokument}
\label{sec:ueber-dokument}
Diese Arbeit wurde in \LaTeX{} verfasst. Diese Art der Dokumentation bietet gegenüber den normalen Textverarbeitungen gewisse Vorteile hinsichtlich der Formatierung und des Einbindens von Grafiken. Auch Formeln können sehr einfach und effizient angegeben werden.

\clearpage

\newpage
\thispagestyle{empty}
\mbox{}

\clearpage

\section*{Projektteam}
\label{sec:projektteam}

\subsection*{Lukas Freyler}
\begin{wrapfigure}[10]{l}{0.5\textwidth}
\begin{center}
  \includegraphics[width=0.35\textwidth]{fig/logoMecha}
\end{center}
\end{wrapfigure}
\mbox{}\\
\mbox{}\\
\textbf{Aufgabenbereich}:\\
\LaTeX{}\\
\textbf{Betreuer}:\\
Dipl.-Ing. Manfred Steiner
\mbox{}\\
\mbox{}\\
\mbox{}\\
\mbox{}\\
\mbox{}\\
\mbox{}\\

\subsection*{Stefan Ornik}
\begin{wrapfigure}[10]{l}{0.5\textwidth}
\begin{center}
  \includegraphics[width=0.35\textwidth]{fig/logoMecha}
\end{center}
\end{wrapfigure}
\mbox{}\\
\mbox{}\\
\textbf{Aufgabenbereich}:\\
\LaTeX{}\\
\textbf{Betreuer}:\\
Dipl.-Ing. Wolfgang Mader
\mbox{}\\
\mbox{}\\
\mbox{}\\
\mbox{}\\
\mbox{}\\
\newpage

\subsection*{Fabio Pölzl}
\begin{wrapfigure}[10]{l}{0.5\textwidth}
\begin{center}
  \includegraphics[width=0.35\textwidth]{fig/logoMecha}
\end{center}
\end{wrapfigure}
\mbox{}\\
\mbox{}\\
\textbf{Aufgabenbereich}:\\
\LaTeX{}\\
\textbf{Betreuer}:\\
Dipl.-Ing. Werner Harnisch
\mbox{}\\
\mbox{}\\
\mbox{}\\
\mbox{}\\
\mbox{}\\
\mbox{}\\

\subsection*{Dominik Riegelnegg}
\begin{wrapfigure}[10]{l}{0.5\textwidth}
\begin{center}
  \includegraphics[width=0.35\textwidth]{fig/logoMecha}
\end{center}
\end{wrapfigure}
\mbox{}\\
\mbox{}\\
\textbf{Aufgabenbereich}:\\
\LaTeX{}\\
\textbf{Betreuer}:\\
BEd. Otto Schuller
\mbox{}\\
\mbox{}\\
\mbox{}\\
\mbox{}\\
\mbox{}\\
\newpage